\section{Do official values respond to a flood? A pre-specified design}\label{sec:design}

The identifying variation for Question~1 is the fixed calendar of valuation rounds. The 2023--26 round was set before the 2024 northern floods and the 2025 floods. The first round published after those floods will reveal whether the committees revised the values of flooded parcels relative to nearby unflooded ones. This section describes the design, the flood maps and their validation, the samples, the balance and power evidence, and the Monte Carlo and synthetic checks. The design is written down in a pre-analysis plan with the treatment and control lists hashed; registration on OSF is pending and is to precede the announcement of the next round, so the design is pre-specified but not yet pre-registered. The primary estimate does not exist yet; Section~\ref{sec:status} explains why and when. The plan carries the registration detail (fallback orders, file names, hash rules, stopping rules); this section keeps what a reader needs to judge identification.

\subsection{Treatment, controls and estimator}

The unit is the cadastral sheet. The outcome is the change in the sheet's mean log value per square wah between the 2023--26 round and the next round, computed over land units that exist in both rounds (``pseudo land units''), so that re-zoning does not enter the outcome. A sheet is \emph{treated} if more than 25\% of its land pixels were flooded in the 2024 radar maps and at most 25\% in 2025, so that the treated sheets are those flooded once, in the event under study. \emph{Controls} are sheets within 5~km of a treated sheet, in the same province, with less than 1\% flooded in every source (2024 radar, 2025 radar, UNOSAT, and Sentinel-2 optical), and with height above the valley floor no greater than the 90th percentile of treated sheets, so that controls lie in the same floodplain rather than on the surrounding hills. Provinces enter the primary sample if they have at least 25 treated sheets. Four provinces qualify in 2024: Chiang Rai (\radarTreatedChiangRai{} treated), Nong Khai (\radarTreatedNongKhai), Nakhon Phanom (\radarTreatedNakhonPhanom) and Phayao (\radarTreatedPhayao), giving \radarTreated{} treated and \radarControls{} control sheets in \radarClusters{} 1:50,000 clusters, \radarTreatedClusters{} of which contain treated sheets (Appendix Table~\ref{tab:radarprov}). A second 2024 sample defines treatment from UNOSAT's flood extents instead, with the same 25\% threshold on the share of land pixels inside the extent, and a 2025 sample applies the radar rule to the 2025 floods on sheets that were dry in 2024 (Section~\ref{sec:samples}).

Within each province, controls are re-weighted by entropy balancing to match the treated sheets' means of the 2023--26 log value, height above valley floor, built-up share, cropland share and a rural indicator; if the full set cannot be balanced the worst-fitting variable is dropped until every standardised gap is at most 0.01. The estimator is weighted least squares of the value change on the treatment indicator with province fixed effects, elevation, distance to permanent water, land-cover shares and the rural indicator as covariates. Inference is by randomisation: 1,000 placebo flood footprints are drawn by shifting the real 2024 footprint within the province and re-selecting treated and control sheets, and the primary $p$-value is the share of placebo estimates at least as large in absolute value as the actual one. A wild cluster bootstrap (Rademacher weights, 9,999 draws, restricted, CR1 $t$-statistic, clusters at the 1:50,000 block) is the check, and is primary for the UNOSAT sample, whose footprint cannot be shifted in the same way. The two 2024 samples form one Holm family; the 2025 sample is a separate family. The pre-specified decision rule is a negative coefficient with a Holm-adjusted $p < 0.05$ in at least one 2024 sample. Secondary analyses are the interaction of treatment with height above valley floor, a descriptive carry-over share (the share of treated sheets whose land units are re-drawn between rounds, with Lee bounds on the moved-unit change), and the sensitivities of Appendix~\ref{app:sens}.

The coefficient is a reduced-form quantity: the response of official values, in the first round after the event, to a realised flood. Equation~\eqref{eq:valuer} has no term for realised damage, so that response combines physical damage, any updating of the hazard the committees attach to the land, administrative or relief-driven revision, and changes in the comparables the committees use. Reading it as $\lambda_T\delta$ requires an outside value of the discount $\delta$, which Thailand's lack of public transaction prices does not supply; Section~\ref{sec:status} returns to what the estimate can and cannot say about $\lambda_T$.

\subsection{Validating the radar maps}\label{sec:validation}

Radar flood maps can err in two ways that matter for the design: they can miss flooded sheets, which then enter the control group and attenuate the estimate, and they can flag dry sheets as flooded, which dilutes the treated group. I measure both against Sentinel-2 optical imagery, which is independent of the radar and accurate where cloud permits. For each sheet and each scene with at least half its pixels clear, the water share is the fraction of clear pixels classified as water in the scene classification layer; ``new water'' is the largest 2024 water share minus the largest 2023 water share over the same window, and a sheet is optically flooded if new water exceeds 10\%, a lower bar than the radar's 25\% because clear optical scenes are rare and cloud hides the flood peak. Optical water classification misses water under vegetation and among buildings, so for treated sheets the optical rate is a lower bound on precision.

\begin{table}[t]
\centering\footnotesize
\caption{Radar flood maps against Sentinel-2 optical imagery, 2024 primary sample}\label{tab:validation}
% Generated by 04_analysis/src/paper_numbers.py. Do not edit.
\begin{adjustbox}{max width=\textwidth}
\begin{tabular}{lrrrrr}
\toprule
 & \multicolumn{2}{c}{Controls} & \multicolumn{3}{c}{Treated} \\
\cmidrule(lr){2-3}\cmidrule(lr){4-6}
Province & Sheets & Optically flooded (\%) & Sheets & Optically flooded (\%) & Median new-water share (\%) \\
\midrule
Chiang Rai & 495 & 0.6 & 88 & 56.8 & 14.3 \\
Nakhon Phanom & 981 & 0.3 & 25 & 88.0 & 31.1 \\
Nong Khai & 333 & 5.4 & 60 & 78.3 & 46.5 \\
Phayao & 82 & 0.0 & 25 & 64.0 & 28.2 \\
\midrule
Pooled & 1,891 & 1.3 & 198 & 68.2 & 24.6 \\
\bottomrule
\end{tabular}
\end{adjustbox}

\tabnote{A sheet is ``optically flooded'' if Sentinel-2 shows new open water on it in the flood period. Treated sheets: more than 25\% flooded in the Sentinel-1 radar map; controls: less than 1\% flooded in every source, within 5~km of a treated sheet, height-capped. The median new-water share is over treated sheets. Source: \texttt{s1\_s2\_validation.csv}.}
\end{table}

Table~\ref{tab:validation} gives the result. Of \radarControls{} control sheets, \valControlContam\% (\valControlContamN{} sheets) show optical flooding, \valNongKhaiContamN{} of them in Nong Khai (\valNongKhaiContam\% of its controls) and the rest in Chiang Rai and Nakhon Phanom; in Phayao none does. Of \radarTreated{} treated sheets, \valTreatedPrec\% show optical flooding, with a median new-water share of \valTreatedNewWater\%; precision ranges from \valChiangRaiPrec\% in Chiang Rai to \valNakhonPhanomPrec\% in Nakhon Phanom. The treated group therefore contains some sheets that the optical sensor did not see flooded, a mixture of genuine radar false positives and optical misses under cloud and vegetation. If this misclassification is unrelated to what the committees know about the land, it attenuates the estimate towards zero and does not create a spurious effect. The qualification matters: if the false positives are seasonally wet paddy whose hazard the committees already price, the error is not benign, which is why the treated sheets' land cover and the reference-season rule are examined below.

An independent check uses UNOSAT's event analyses of the 2024 floods, which cover Chiang Rai and Phayao. On the \unoSheets{} sheets in their extent, UNOSAT flags \unoTreated{} sheets as flooded and the radar maps \unoRadarTreated, with \unoBoth{} in common: precision of the radar classification against UNOSAT is \unoPrecision{} but recall only \unoRecall, agreement \unoAgreement{} and Cohen's kappa \unoKappa, with a correlation of \unoCorr{} between the two flood shares. The radar maps, built from a few scenes per orbit, see a small subset of what UNOSAT's multi-date analysis sees. This is why the design carries a second 2024 sample in which UNOSAT defines treatment, and why the radar sample's controls are required to be dry in every source including UNOSAT.

\paragraph{What the treated sheets are.} The radar rule counts a pixel as flooded only if it was water in no 2023 scene of the same window, which is meant to remove paddy that floods every year. A single reference season cannot do so where 2023 happened to be dry, and the variant that also screens out water seen in 2021 and 2022 removes \sensRefDropped{} of the \radarTreated{} treated sheets (\sensRefDroppedPct\%; Appendix Table~\ref{tab:sensitivities}), which is direct evidence that the single-season rule admits some recurrent water. Table~\ref{tab:treatedlc} shows what the treated sheets are: cropland is the largest 2021 land-cover class on \tlcCropDominant{} of the \tlcTreated{} treated sheets and exceeds half the sheet on \tlcCropOverHalf; built-up land is the largest class on \tlcBuiltDominant{} and exceeds a fifth of the sheet on \tlcBuiltOverTwenty. The treated units are, overwhelmingly, northern and north-eastern rice land. Two consequences follow. The design's external validity for bank collateral, which is largely residential and urban, is limited, and the chain from $\lambda_T$ to $\lambda_B$ in Section~\ref{sec:stress} runs through a population that this sample cannot speak to directly. And the cropland-excluded variant keeps only \sensNocropTreated{} treated sheets in \sensNocropClusters{} clusters, so no separate power claim is made for built-up sheets. Whether the multi-year reference period should become the primary treatment definition is a decision to be taken before registration; the plan currently keeps the single-season rule as primary and the 2021--23 rule as a registered sensitivity.

\begin{table}[t]
\centering\footnotesize
\caption{Land cover of the treated sheets, 2024 radar sample}\label{tab:treatedlc}
% Generated by 04_analysis/src/paper_numbers.py. Do not edit.
\begin{adjustbox}{max width=\textwidth}
\begin{tabular}{lrrrr}
\toprule
 & \multicolumn{3}{c}{Treated sheets} & Controls \\
\cmidrule(lr){2-4}\cmidrule(lr){5-5}
WorldCover 2021 class & Largest share & Share above 50\% & Mean share (\%) & Mean share (\%) \\
\midrule
Built-up & 4 & 4 & 4.5 & 14.0 \\
Cropland & 171 & 158 & 65.8 & 53.0 \\
Tree cover & 16 & 8 & 19.3 & 27.2 \\
Permanent water & 7 & 3 & 5.1 & 1.4 \\
\bottomrule
\end{tabular}
\end{adjustbox}

\tabnote{WorldCover 2021 shares of each sheet's 10~m pixels. ``Largest share'': number of treated sheets on which the class has the largest share of the four; ``share above 50\%'': number on which the class covers more than half the sheet. Controls: the \radarControls{} primary controls, unweighted. Source: \texttt{rq1\_2024\_sample\_sheets.csv.gz}.}
\end{table}

\subsection{Balance and power}\label{sec:balance}

\begin{table}[t]
\centering\footnotesize
\caption{Balance of treated and control sheets, 2024 primary sample}\label{tab:balance}
% Generated by 04_analysis/src/paper_numbers.py. Do not edit.
\begin{adjustbox}{max width=\textwidth}
\begin{tabular}{lcc}
\toprule
 & \multicolumn{2}{c}{Standardised mean difference, treated minus controls} \\
\cmidrule(lr){2-3}
Covariate & Unweighted & Registered weights \\
\midrule
Log value per sq.\ wah, 2023--26 & $-$0.41 & 0.08 \\
Height above valley floor (m) & $-$0.15 & 0.32 \\
Built-up share & $-$0.68 & $-$0.00 \\
Cropland share & 0.52 & 0.13 \\
Tree-cover share & $-$0.48 & $-$0.38 \\
Permanent-water share & 0.42 & 0.67 \\
\midrule
Effective number of controls & 1,891 & 98.0 \\
MDE (SD 0.10, ICC 0.2) & 0.047 & 0.053 \\
\bottomrule
\end{tabular}
\end{adjustbox}

\tabnote{Standardised mean differences (treated minus controls, divided by the pooled standard deviation) before weighting and after entropy balancing with the registered weights (within province, on log value, height, built-up and cropland shares and the rural indicator, with the fallback order of the plan; pooled row of Appendix Table~\ref{tab:regbalance}). Tree and water shares are not balanced on and are reported as a check. The effective number of controls is the Kish effective sample size of the weights. MDE: minimum detectable effect at 80\% power and 5\% size, two-sided, for a standard deviation of value changes of 0.10 and intra-cluster correlation 0.2, with $t$ critical values on the number of clusters minus one. Sources: \texttt{rq1\_2024\_pooled\_balance.csv} (unweighted), \texttt{rq1\_round\_change\_balance.csv} (registered weights), \texttt{rq1\_2024\_pooled\_mde.csv}.}
\end{table}

Flooded sheets are not a random sample of the floodplain. Before weighting, treated sheets have lower values (standardised difference \radarSmdValue), less built-up land (\radarSmdBuilt) and more cropland (\radarSmdCrop) than controls that are already restricted to the same floodplain (Table~\ref{tab:balance}). Under the registered weights the pooled gaps on the balancing variables fall to \radarGapValueW{} on value, \radarGapHeightW{} on height, \radarGapBuiltW{} on built-up share and \radarGapCropW{} on cropland, and the tree-cover and water shares, which are not balanced on, stand at \radarSmdTreeW{} and \radarSmdWaterW. The gaps are small but not zero, because balancing is done province by province and the pooled gap aggregates provinces with different fallbacks. The price is a loss of effective sample: the \radarControls{} controls carry the weight of \radarEssBalance{} after balancing. Appendix Table~\ref{tab:regbalance} gives the province detail, and it is uneven. In Phayao, \regTreatedPhayao{} treated sheets against \regControlsPhayao{} controls leave \regEssPhayao{} effective controls, the registered fallback keeps only the built-up share and the rural indicator, and the remaining gaps are \regGapValuePhayao{} standard deviations on value, \regGapHeightPhayao{} on height, \regGapCropPhayao{} on cropland, \regGapTreePhayao{} on tree cover and \regGapWaterPhayao{} on water. In Nakhon Phanom \regControlsNakhonPhanom{} controls shrink to \regEssNakhonPhanom{} effective controls. Gaps of this size mean that Phayao's treated sheets lie partly outside the support of its controls, so the province contributes extrapolation rather than comparison. Two responses are added to the plan before registration: a sensitivity that trims treated and control sheets to common support on height and water share, and province-by-province estimates with their effective sample sizes reported beside the pooled one. Whether Phayao stays in the primary sample under an overlap rule is a decision to be fixed in the plan before any round is published, not after.

Power is computed from the cluster structure rather than from the sheet count. With $G$ clusters and treated share $a_g$ and control share $b_g$ of the weights in cluster $g$, the variance of the difference in means is $\sigma^2[\rho\sum_g (a_g - b_g)^2 + (1-\rho)(1/n_1 + 1/n_0)]$ and the minimum detectable effect is $(t_{0.975,G-1} + t_{0.80,G-1})$ times its square root. With $\sigma = 0.10$ and $\rho = 0.2$ the registered primary sample has a minimum detectable effect of \regMdePooled{} log points after weighting (\radarMde{} unweighted); the unweighted figure is \radarMdeIccLow{} at $\rho = 0.05$ and \radarMdeIccHigh{} at $\rho = 0.5$, and \radarMdeHigh{} if value changes are twice as dispersed. The equivalence margin that the unweighted design could rule out at 80\% power is \radarEquiv. For the carry-over share the minimum detectable difference is \radarCarry{} percentage points. The scale of $\sigma$ is uncertain because no earlier public round exists; the plan treats 0.10 as the central case and 0.20 as the pessimistic one.

\begin{table}[t]
\centering\footnotesize
\caption{Power against the effects the framework implies}\label{tab:powerlambda}
% Generated by 04_analysis/src/paper_numbers.py. Do not edit.
\begin{adjustbox}{max width=\textwidth}
\begin{tabular}{llrrrr}
\toprule
$\delta$ & $\lambda_T$ & Sheet effect & \multicolumn{3}{c}{Power at the registered MDE} \\
\cmidrule(lr){4-6}
 & & & Radar 2024 & UNOSAT 2024 & Radar 2025 \\
\midrule
2.3 & 0.25 & $-$0.004 & 0.05 & 0.05 & 0.07 \\
2.3 & 0.50 & $-$0.008 & 0.07 & 0.06 & 0.11 \\
2.3 & 1.00 & $-$0.016 & 0.13 & 0.11 & 0.32 \\
\midrule
4.6 & 0.25 & $-$0.008 & 0.07 & 0.06 & 0.11 \\
4.6 & 0.50 & $-$0.016 & 0.13 & 0.11 & 0.32 \\
4.6 & 1.00 & $-$0.032 & 0.39 & 0.30 & 0.85 \\
\midrule
9.2 & 0.25 & $-$0.016 & 0.13 & 0.11 & 0.32 \\
9.2 & 0.50 & $-$0.032 & 0.39 & 0.30 & 0.85 \\
9.2 & 1.00 & $-$0.064 & 0.92 & 0.84 & 1.00 \\
\midrule
Damage 10\% & -- & $-$0.070 & 0.96 & 0.89 & 1.00 \\
Damage 20\% & -- & $-$0.140 & 1.00 & 1.00 & 1.00 \\
\bottomrule
\end{tabular}
\end{adjustbox}

\tabnote{Sheet effect: $0.7 \times \lambda_T \times \delta$, the first-round fall in the sheet mean when the committees capture a share $\lambda_T$ of a full-information discount $\delta$ (in per cent) and 30\% of land units carry their value over unchanged, as in the synthetic rounds of Table~\ref{tab:montecarlo}. ``Damage'' rows: a realised-damage discount of 10\% or 20\%, fully passed through. Power: two-sided test at 5\% from the registered weighted MDE of each sample (Table~\ref{tab:samples}), inverting $\text{MDE} = (t_{0.975,G-1} + t_{0.80,G-1}) \times \text{SE}$. Source: \texttt{rq1\_round\_change\_balance.csv}, \texttt{rq1\_round\_change\_unosat\_balance.csv}, \texttt{rq1\_round\_change\_2025\_balance.csv}; author's calculations.}
\end{table}

What a minimum detectable effect of \regMdePooled{} means for the framework's parameters is shown in Table~\ref{tab:powerlambda}. With $\delta = 4.6\%$, the inland-floodplain discount of \citet{beltran2018flood}, full recognition ($\lambda_T = 1$) implies a sheet effect of \effFullMid{} once 30\% of units carry over, and the primary sample's power against it is \powFullMid; against half recognition (\effHalfMid) it is \powHalfMid. Power is adequate only for a discount twice that size (\powFullHigh{} at $\delta = 9.2\%$) or for realised-damage discounts of the size reported after Hurricane Sandy (\powDamageTen{} at 10\%, \powDamageTwenty{} at 20\%; \citealp{ortega2018rising,guo2026rising}). The 2025 sample does better (\powFullMidEv{} against full recognition of 4.6\%, \powHalfMidEv{} against half), which is one reason the plan registers a stacked 2024-plus-2025 estimate as a sensitivity. In the same units, the registered equivalence margin of \radarEquiv{} corresponds to $\lambda_T \approx \equivLambdaMid$ at $\delta = 4.6\%$ and $\lambda_T \approx \equivLambdaHigh$ at $9.2\%$, or to a fully passed-through realised-damage discount of about \equivDamage\%: a null would rule out responses larger than a fully recognised hazard discount, not partial recognition.

\subsection{The three samples}\label{sec:samples}

\begin{table}[t]
\centering\footnotesize
\caption{The three registered samples}\label{tab:samples}
% Generated by 04_analysis/src/paper_numbers.py. Do not edit.
\begin{adjustbox}{max width=\textwidth}
\begin{tabular}{lccc}
\toprule
 & Radar 2024 (primary) & UNOSAT 2024 & Radar 2025 \\
\midrule
Treated sheets & 198 & 1,214 & 437 \\
Control sheets & 1,891 & 579 & 7,664 \\
Effective controls after weighting & 98.0 & 103.9 & 633.6 \\
Clusters (1:50,000 sheets) & 57 & 28 & 134 \\
Treated clusters & 21 & 18 & 49 \\
Provinces & 4 & 2 & 9 \\
MDE, unweighted (SD 0.10, ICC 0.2) & 0.047 & 0.065 & 0.029 \\
MDE, weighted & 0.053 & 0.061 & 0.030 \\
Carry-over MDD (pp) & 23.5 & 32.3 & 14.6 \\
\bottomrule
\end{tabular}
\end{adjustbox}

\tabnote{Radar 2024: treated if more than 25\% flooded in Sentinel-1 maps of the 2024 northern floods and at most 25\% in 2025, four provinces. UNOSAT 2024: treated if more than 25\% of land pixels fall inside UNOSAT's flood extents, Chiang Rai and Phayao. Radar 2025: treated if more than 25\% flooded in 2025 radar maps and under 1\% in 2024, nine provinces. MDE at 80\% power, $\sigma = 0.10$, ICC 0.2. Carry-over MDD: minimum detectable difference in the share of treated sheets whose land units are re-drawn. Sources: \texttt{rq1\_2024\_pooled\_mde.csv}, \texttt{rq1\_2024\_mde\_changes.csv}, \texttt{rq1\_sample\_sensitivity.csv}, \texttt{rq1\_round\_change\_*balance.csv}.}
\end{table}

Table~\ref{tab:samples} places the primary sample beside the two others. The UNOSAT sample is larger in treated sheets (\unoTreatedReg) but thin in controls (\unoControls{} after dropping \unoRadarWet{} comparable sheets that are wet in the radar maps), with \unoTreatedClusters{} treated clusters and a minimum detectable effect of \unoMdeW{} after weighting; its ex-ante balance is in Appendix Table~\ref{tab:unosatbalance}, and its weakness is Phayao, which holds \unoPhayaoTreated{} of the treated sheets but only \unoPhayaoControls{} controls. The 2025 sample is the largest: \evTreated{} treated and \evControls{} control sheets in \evProvinces{} provinces (\evProvinceList) and \evClusters{} clusters, with a minimum detectable effect of \evMdeW{} and a carry-over MDD of \evCarry{} points. Its balance is poor on two variables. In Phatthalung the height variable had to be dropped, leaving a standardised gap of \evPhatthalungGapHeight{} on it; and the gaps on the water share, which is not balanced on, are \evPhatthalungGapWater{} in Phatthalung, \evPattaniGapWater{} in Pattani, \evNakhonSiGapWater{} in Nakhon Si Thammarat and \evSongkhlaGapWater{} in Songkhla (\evGapWater{} pooled). Gaps of several standard deviations mean that the treated riverside and lagoon sheets there lie outside the controls' support, so the common-support trimming described above applies to this sample with more force than to the primary one. The threshold of 25 treated sheets per province was chosen before the samples were built; Appendix Table~\ref{tab:thresholds} shows that 20 would add Nakhon Si Thammarat and lower the 2024 minimum detectable effect to 0.044, while 30 would leave two provinces and raise it to 0.060.

\subsection{Monte Carlo and synthetic checks}\label{sec:montecarlo}

Power formulas assume the estimator behaves. To check that it does, I generate synthetic next rounds from the real 2023--26 values: each land unit's value is multiplied by a province-level revaluation factor and idiosyncratic sheet and unit noise, a planted treatment effect is applied to the land units of treated sheets, 30\% of units carry their value over unchanged so that 70\% of the planted effect reaches the sheet mean, and a random 2\% of sheets are made to attrit through re-drawn units. The full registered pipeline, from sample construction through balancing to inference, is run on each replication. Table~\ref{tab:montecarlo} reports \mcRadarReps{} replications for each of the three samples with the wild cluster bootstrap as the test; Table~\ref{tab:mcextended} reports \mcxReps{} replications for the primary sample, adds the registered placebo-footprint test on the first \mcxRiReps{} of them, and records how the secondary analyses behave.

\begin{table}[t]
\centering\footnotesize
\caption{Monte Carlo performance of the registered estimator on synthetic next rounds}\label{tab:montecarlo}
% Generated by 04_analysis/src/paper_numbers.py. Do not edit.
\begin{adjustbox}{max width=\textwidth}
\begin{tabular}{llrrrrr}
\toprule
Sample & Planted & Expected sheet effect & Mean estimate & SD of estimates & Mean cluster SE & Rejection rate (5\%) \\
\midrule
Radar 2024 & 0.00 & 0.000 & $-$0.0006 & 0.0155 & 0.0132 & 0.05 \\
Radar 2024 & $-$0.05 & $-$0.035 & $-$0.0353 & 0.0162 & 0.0127 & 0.63 \\
UNOSAT 2024 & 0.00 & 0.000 & 0.0020 & 0.0158 & 0.0130 & 0.03 \\
UNOSAT 2024 & $-$0.05 & $-$0.035 & $-$0.0332 & 0.0175 & 0.0130 & 0.53 \\
Radar 2025 & 0.00 & 0.000 & $-$0.0007 & 0.0081 & 0.0074 & 0.04 \\
Radar 2025 & $-$0.05 & $-$0.035 & $-$0.0356 & 0.0075 & 0.0072 & 1.00 \\
\bottomrule
\end{tabular}
\end{adjustbox}

\tabnote{\mcRadarReps{} replications per cell. Planted: the log change applied to every land unit on a treated sheet; the expected sheet effect is 70\% of it. Rejection rate: share of replications with a wild cluster bootstrap $p < 0.05$; its Monte Carlo standard error at a true rate of 5\% is \mcSeHundred{} percentage points. Synthetic next rounds are built from the real 2023--26 values with province-level revaluation and sheet and unit noise; they are not data. Sources: \texttt{rq1\_round\_change\_montecarlo.csv}, \texttt{rq1\_round\_change\_unosat\_montecarlo.csv}, \texttt{rq1\_round\_change\_2025\_montecarlo.csv}.}
\end{table}

\begin{table}[t]
\centering\footnotesize
\caption{Extended Monte Carlo, 2024 primary sample: both tests and the secondary analyses}\label{tab:mcextended}
% Generated by 04_analysis/src/paper_numbers.py. Do not edit.
\begin{adjustbox}{max width=\textwidth}
\begin{tabular}{llrrrrrrr}
\toprule
Planted & Expected & Mean estimate & SD & Mean SE & \multicolumn{2}{c}{Rejection rate (5\%)} & $|t|$ of H1a $> 1.96$ & Lee bounds exclude 0 \\
\cmidrule(lr){6-7}
 & & & & & Bootstrap & Placebo & & \\
\midrule
0.00 & 0.000 & 0.0001 & 0.0152 & 0.0129 & 0.055 & 0.195 & 0.085 & 0.547 \\
$-$0.05 & $-$0.035 & $-$0.0350 & 0.0155 & 0.0127 & 0.597 & 0.795 & 0.111 & 0.972 \\
\bottomrule
\end{tabular}
\end{adjustbox}

\tabnote{\mcxReps{} synthetic next rounds per planted effect, same generator as Table~\ref{tab:montecarlo}. Bootstrap: wild cluster bootstrap, 999 draws. Placebo: the registered placebo-footprint test, 500 shifted footprints, on the first \mcxRiReps{} replications (Monte Carlo standard error \mcxRiSeNull{} points at 5\%). $|t|$ of H1a: share of replications in which the treatment-by-height interaction, which has no planted heterogeneity, has a cluster $t$-statistic above 1.96. Lee bounds exclude 0: share in which the point bounds on the moved-unit change lie wholly on one side of zero. The placebo test does not hold its size under the null; the text explains why. Not results. Source: \texttt{rq1\_montecarlo\_extended.csv}.}
\end{table}

Under a true null the primary 2024 estimator rejects \mcRadarSize{} of the time in Table~\ref{tab:montecarlo}, the UNOSAT estimator \mcUnoSize{} and the 2025 estimator \mcEvSize, each within the Monte Carlo error of \mcSeHundred{} points around 5\%; in the \mcxReps{} replications of Table~\ref{tab:mcextended} the bootstrap rejects \mcxNullWcb{} of the time, but the placebo-footprint test rejects \mcxNullRi{} (on \mcxRiReps{} replications), nearly four times its nominal size, a failure examined below. With a planted $-0.05$ per land unit (expected sheet effect \mcExpected) the mean estimate is \mcxEffMean{} with a standard deviation across replications of \mcxEffSd{} against a mean cluster standard error of \mcxEffSe: the estimator is close to unbiased, but the CR1 standard error understates the sampling standard deviation by about a fifth (\mcxEffSeSdRatio\% of it), which is why inference rests on the bootstrap and the placebo footprints rather than on the analytic standard error. Power against that effect is \mcxEffWcb{} with the bootstrap and \mcxEffRi{} with the placebo test; it is \mcUnoPower{} for UNOSAT and \mcEvPower{} for 2025 in Table~\ref{tab:montecarlo}. The Monte Carlo power is higher than the analytic figure in Table~\ref{tab:powerlambda} because the synthetic rounds' sampling standard deviation (\mcxEffSd) is below the \regMdePooled-based standard error that the formula assumes; the analytic figure is the registered one.

The secondary analyses behave as follows under the null. The height interaction, with no planted heterogeneity, has $|t| > 1.96$ in \mcxNullInterT\% of replications, so a single ``significant'' interaction on real data, like the one on the null draw of Appendix Table~\ref{tab:synspecs}, would be unremarkable, and the plan treats H1a as secondary. The Lee point bounds on the moved-unit change exclude zero in \mcxNullLee\% of null replications (mean bounds [\mcxNullLeeLower, \mcxNullLeeUpper] around a mean naive estimate of \mcxNullLeeNaive): the moved units are, by construction of the synthetic rounds, those whose value was re-drawn, so their change is not centred on zero even without a treatment effect, and the bounds are a description of attrition rather than a test. Inference with \radarTreatedClusters{} treated clusters in four provinces is a known difficulty for the wild cluster bootstrap \citep{mackinnon2017wild,mackinnon2018few}; it is why the draft plan made the placebo-footprint test primary and the bootstrap a check, and why both tests are simulated here rather than assumed. The simulation reverses that ranking.

The placebo-footprint test over-rejects because its reference distribution is too tight. A diagnostic on \pdxDraws{} synthetic null rounds with \pdxFootprints{} shifted footprints each (\texttt{rq1\_placebo\_diagnostic.py}) compares the placebo estimates with the sampling distribution of the actual estimate. The standard deviation of the placebo estimates is \pdxPlaceboSdLo{} to \pdxPlaceboSdHi{} across the three rounds, about \pdxSdRatio\% of the standard deviation of the actual estimate under the null across replications (\mcxNullSd), so an actual estimate drawn from the null falls in the tails of the placebo distribution far more often than 5\% of the time. The shifted footprints cover as many 1:50,000 clusters as the real one (\pdxPlaceboClusters{} on average against \pdxActualClusters) but rest on many more effective controls: after balancing, the Kish effective number of controls behind a placebo estimate is about \pdxPlaceboEss, against \pdxActualEss{} for the treated sheets (Table~\ref{tab:regbalance}). The real treated sheets, low riverside rice land, lie at the edge of the controls' support on height and land cover, so balancing concentrates the weight on a few controls; a shifted footprint, which the rule places only on eligible dry controls, lies in the middle of that support and is estimated with far more precision. The test's exchangeability assumption, that the actual footprint is one draw from the same population as the placebo footprints, fails on precision rather than on location, and the placebo mean is close to zero in every round. The bootstrap, which resamples the actual footprint's own clusters, has no such problem and holds its size at \radarTreatedClusters{} treated clusters. Two repairs were then tested on the same synthetic rounds (Table~\ref{tab:placeborepairs}; \texttt{rq1\_placebo\_restricted\_mc.py}, \prxReps{} replications per planted effect, 200 footprints per rule). The first restricts the pool: a logit of treatment on the balancing variables with province effects is fitted, and footprints may land only on controls whose fitted propensity is at least the 10th percentile of the treated sheets', with a stricter variant that starts at the median. It does not work. The effective controls behind a placebo estimate fall from \prxRegEss{} to \prxResEss{} and \prxStrEss, but the footprints, snapping to a sparser pool, scatter over \prxResClusters{} and \prxStrClusters{} clusters against \prxRegClusters, and the two effects cancel: the placebo standard deviation stays at about \prxResSd{} and the size at \prxNullRes{} and \prxNullStr. The second repair studentises the comparison, ranking $|\hat\beta / \widehat{\mathrm{se}}|$ rather than $|\hat\beta|$; the failure it addresses is the one \citet{chung2013studentised} describe, that a permutation test need not hold its level when the groups differ in more than the parameter under test. With the registered pool its size is \prxNullRegT{} on the same replications on which the bootstrap's is \prxNullWcb{}, and its power against the planted $-0.05$ is \prxEffRegT{} against the bootstrap's \prxEffWcb. It works because the cluster standard error understates the sampling standard deviation by a similar proportion for the real and the placebo footprints (the $t$-statistic has a standard deviation of \prxNullTSd{} across replications and \prxRegTSd{} across placebo draws), so the studentised ranking removes the precision gap that the raw ranking cannot see. The plan is therefore amended before registration in two parts: the wild cluster bootstrap becomes the primary test for the radar samples, and the placebo-footprint check is studentised; the unstudentised test and the restricted pools are dropped. The decision is recorded in Section~\ref{sec:discussion} with the other pre-registration decisions.

\begin{table}[t]
\centering\footnotesize
\caption{Repairs to the placebo-footprint test on synthetic next rounds, 2024 primary sample}\label{tab:placeborepairs}
% Generated by 04_analysis/src/paper_numbers.py. Do not edit.
\begin{adjustbox}{max width=\textwidth}
\begin{tabular}{lrrrrrrrr}
\toprule
Placebo rule & Eff.\ controls & Clusters & SD of estimates & SD of $t$ & \multicolumn{2}{c}{Size (null)} & \multicolumn{2}{c}{Power ($-0.05$)} \\
\cmidrule(lr){6-7}\cmidrule(lr){8-9}
 & & & & & Estimate & $t$ & Estimate & $t$ \\
\midrule
Registered pool & 510 & 20.7 & 0.0110 & 1.22 & 0.195 & 0.055 & 0.80 & 0.56 \\
Restricted pool (25/10/5, $2\times$) & 234 & 25.5 & 0.0109 & 1.11 & 0.200 & 0.080 & 0.79 & 0.60 \\
Strict pool (50/25/10/5, $1.5\times$) & 177 & 27.6 & 0.0109 & 1.09 & 0.200 & 0.075 & 0.80 & 0.64 \\
\midrule
Wild cluster bootstrap & 98 & 21 & 0.0164 & 1.25 & \multicolumn{2}{c}{0.070} & \multicolumn{2}{c}{0.59} \\
\bottomrule
\end{tabular}
\end{adjustbox}

\tabnote{\prxReps{} synthetic next rounds per planted effect, the first \prxReps{} draws of Table~\ref{tab:mcextended}; 200 placebo footprints per rule and replication. Registered pool: every control within the treated height range. Restricted pool: controls whose fitted propensity is at least the 25th, 10th or 5th percentile of the treated sheets' propensities, the largest cut that leaves twice as many eligible as treated sheets in the province (the 10th in every province on these draws). Strict pool: the ladder starts at the median and accepts 1.5 times the treated count. Eff.\ controls: Kish effective number of controls after balancing, mean over footprints; for the bootstrap row, the real footprint's. SD of estimates and of $t$: across placebo draws, mean over replications; for the bootstrap row, across replications of the actual estimate. Size: share of null replications with $p < 0.05$, ranking placebo draws by $|\hat\beta|$ (Estimate) or by $|\hat\beta / \widehat{\mathrm{se}}|$ ($t$); Monte Carlo standard error \mcxRiSeNull{} points at 5\%. Power: the same at a planted $-0.05$ per land unit. Not results. Source: \texttt{rq1\_placebo\_restricted\_mc.csv}.}
\end{table}

A single synthetic draw run through every registered specification (Appendix Table~\ref{tab:synspecs}) shows what a result would look like: the primary estimate is \synPrimary{} (\synPrimarySe), placebo $p = \synPrimaryRi$, bootstrap $p = \synPrimaryWcb$; the unweighted estimate is closer to the planted value, \synUnweighted{} (\synUnweightedSe), because the planted effect is homogeneous and weighting only costs precision; and attrition is \synAttritionT\% of treated and \synAttritionC\% of control sheets. On that same draw the stacked 2024-plus-2025 estimate is \sensStacked{} (\sensStackedSe) on \sensStackedTreated{} treated sheets in \sensStackedClusters{} clusters, which is the precision the design reaches once both events are in hand. The sensitivities in Appendix Table~\ref{tab:sensitivities} report estimation counts, which are below the sample-file counts because sheets whose change cannot be computed (no land unit present in both rounds) drop out: the no-cap variant keeps \sensNocapControlsEst{} of \sensNocapControls{} controls, the 10\% threshold \sensThrTenTreatedEst{} of \sensThrTenTreated{} treated and \sensThrTenControlsEst{} of \sensThrTenControls{} controls, and the 20-treated province threshold \sensProvTwentyControlsEst{} of \sensProvTwentyControls{} controls. Three rows deserve a note because they are the ones most likely to be wrongly read as results. Dropping sheets with more than 50\% cropland leaves \sensNocropTreated{} treated sheets and \sensNocrop{} (\sensNocropSe). The declared-district and declared-tambon rows are identical because both use a stand-in list, every district holding a treated sheet, in place of the Department of Disaster Prevention and Mitigation's declarations, which have not yet been obtained; the rows show that the code runs, not what the declarations would do. And the row that drops the \sensRefloodDropped{} treated sheets flooded again in 2026 gives \sensReflood{} (\sensRefloodSe) on this draw; Section~\ref{sec:status} explains why that row matters.

\subsection{Status of the primary estimate, timing and what it can show}\label{sec:status}

The design needs the next valuation round. Under the Act, the 2027--30 round would take effect on 1~January 2027 and be published at least thirty days earlier \citep{th2019valuationact}. As of 3--4 November 2026, the Treasury Department's provincial news items record that committees in Chainat, Chachoengsao, Phatthalung and Trat met to consider a one-year extension of the 2023--26 schedule, which the Act allows, the Chainat committee stating an extension to 31~December 2027 \citep{treasury2026chainat,treasury2026chachoengsao,treasury2026phatthalung,treasury2026trat}; provincial offices had been briefed on an extension in late October \citep{treasury2026extbriefing}; no national decision or Gazette notice had been published. The treatment and control files were hashed on 2 November 2026 (Appendix Table~\ref{tab:hashes}); registration is to precede any announcement, and the hashes make a later change to the samples detectable.

Timing enters in two ways, and the plan's decision tree (Appendix~\ref{app:timing}, Table~\ref{tab:timing}) sets them out. The first is when the next round is published: on time, after a one-year extension, or province by province. The second is when its comparables were collected. Committees value by market comparison from registered transactions, so a round published after a flood can reflect it only if its evidence window extends past the flood. The public record does not state the transaction cut-off dates of the next round. The plan therefore carries the case in which a province's survey predates the flood: that province's first round is then a placebo, its primary test moves to the second round, and the case is set province by province from the published survey periods, with the on-time case assumed and the placebo case reported as a sensitivity where the period is unknown. For the 2025 floods, which came later, the placebo case is the default expectation for the first round; the 2025 sample's precision in Table~\ref{tab:samples} is therefore a statement about the second round after the floods unless the survey periods show otherwise. Two further points are pre-specified. Extension is itself decided province by province and may be related to flooding, so the plan records, as a descriptive outcome, whether provinces holding treated sheets extend more often than others; if they do, the set of provinces that revalue on time is selected on treatment and the first-round estimate is reported with that caveat. And under a delayed round the 2026 floods enter the same change: \sensRefloodDropped{} treated sheets, most of them in Nakhon Phanom, flooded again in 2026, so a province whose survey postdates them mixes two floods. The plan keeps those sheets in the primary sample, because re-flooding is an outcome of the same hazard and conditioning on it would select on the cause of the outcome, and reports the 2026-drop row of Appendix Table~\ref{tab:sensitivities} beside the primary estimate.

What the design can deliver, once the round is published, is an estimate of the first-round response of official values to a flood observed on \radarTreated{} to \evTreated{} sheets per event, with a minimum detectable effect of \regMdePooled{} log points in the primary sample, \unoMdeW{} in the UNOSAT sample and \evMdeW{} in the 2025 sample. The benchmarks are of two kinds. Hazard-zone discounts in the advanced-economy literature are 4.6\% for the inland 100-year floodplain \citep{beltran2018flood}, 5.7--8.8\% after hurricanes \citep{bin2013changes} and about 8\% for undamaged homes in New York's flood zone \citep{ortega2018rising}; realised-damage discounts are larger, about 10\% for the sale prices of inundated New York homes \citep{guo2026rising} and 17--22\% for damaged ones \citep{ortega2018rising}. The design is powered against the second kind and against a fully recognised hazard discount at the upper end of the first, not against partial recognition (Table~\ref{tab:powerlambda}). A null at that precision therefore cannot distinguish partial recognition from none, and would not contradict prediction P2, which expects a partial first-round response; the registered equivalence margin would rule out only responses larger than a fully recognised 4.6\% discount. A negative estimate would measure the first-round response to realised flooding, which combines damage, updated hazard and administrative revision; with the within-round trigger of 15\%, it would also say how far the official value had drifted from the market before the round. Only with an outside estimate of $\delta$, which the paper does not have, could it be read as $\lambda_T\delta$.
